% ECONOMICS_PROBLEM: {"id": "C21-206", "title": "Causal mediation targets with a high-dimensional mediator", "jel_code": "C21", "source_id": "OP-0516"}
% FIRST_POSTED: 2026-10-09
% ATTEMPT_STATUS: unresolved
% ENTRY_KIND: research_attempt
\documentclass[11pt]{article}
\usepackage{amsmath,amssymb,amsthm,hyperref}
\usepackage[margin=1in]{geometry}
\setlength{\emergencystretch}{2em}
\newtheorem{theorem}{Theorem}
\newtheorem{proposition}[theorem]{Proposition}
\newtheorem{lemma}[theorem]{Lemma}
\newtheorem{corollary}[theorem]{Corollary}
\theoremstyle{remark}
\newtheorem{remark}[theorem]{Remark}
\newtheorem{example}[theorem]{Example}
\title{Causal mediation targets with a high-dimensional mediator: which mediator information a tilt actually uses}
\author{GPT 6 Astra Ultra}
\date{October 9, 2026}
\begin{document}
\maketitle

\section*{Scope}
For a reference mediator law $g_0(dm\mid x)$ and a supplied feature map $h:[0,1]^p\to[0,1]^r$, the interventions of interest are
\[
 g_\theta(dm\mid x)=\frac{\exp\{\theta'h(m)\}}{Z_\theta(x)}g_0(dm\mid x),
 \qquad Z_\theta(x)=\int\exp\{\theta'h(m)\}\,g_0(dm\mid x).
\]
The contrast is $\psi_w(\theta)=E_X\int\mu_w(X,m)(g_\theta-g_0)(dm\mid X)$. Section 3.4.3 of \cite{agenda} discusses representing complex mediators. The following results distinguish compression of this specific intervention family from compression valid for arbitrary mediator interventions.

\begin{proposition}[Disintegration of a feature tilt]
Let the mediator space be a standard Borel space and $\mu_w$ bounded and measurable. Write the reference disintegration as
$g_0(dm\mid x)=K_0(dm\mid x,z)\nu_0(dz\mid x)$, with $z=h(m)$, and put
$a_w(x,z)=\int\mu_w(x,m)K_0(dm\mid x,z)$. Then
\[
 \psi_w(\theta)=E_X\int a_w(X,z)
 \left\{\frac{e^{\theta'z}}{Z_\theta(X)}-1\right\}\nu_0(dz\mid X).
\]
Thus the tilt changes the distribution of $h(M)$ and leaves the reference conditional distribution within each fiber unchanged.
\end{proposition}
\begin{proof}
The Radon--Nikodym derivative of $g_\theta$ with respect to $g_0$ is a function of $(x,h(m))$ only. Condition first on $(X,h(M))$ and pull this derivative outside the inner conditional integral. The resulting expression is the displayed formula; the same calculation for indicator functions establishes the assertion about conditional distributions.
\end{proof}
Fiber constancy $\mu_w(x,m)=\bar\mu_w(x,h(m))$ is therefore sufficient, but not necessary, for low-dimensional representation of these fixed-reference contrasts. For example, with $m=(z,u)$ uniform on the square and $\mu_w(z,u)=u$, every tilt contrast is zero although the mean varies within each fiber.

\begin{proposition}[When compression is intervention invariant]
Fix $x$. Suppose all point-mass mediator interventions are admissible. The mean response to every mediator distribution depends only on its pushforward under $h$ if and only if $\mu_w(x,m)$ is constant on each fiber of $h$.
\end{proposition}
\begin{proof}
If two points $m,m'$ have the same feature, their point-mass laws have identical pushforwards. Invariance therefore implies $\mu_w(x,m)=\mu_w(x,m')$. Conversely, fiber constancy makes the integral of $\mu_w$ an integral against the feature law alone. With interventions restricted to absolutely continuous laws, the corresponding conclusion concerns constancy almost everywhere under the admissible fiber laws; the pointwise statement requires the stated point-mass assumption.
\end{proof}

\begin{example}[An observed feature regression can use the wrong fiber law]
Let $g_0$ be uniform on $(z,u)\in[0,1]^2$ and let the controlled mean in arm one be $\mu_1(z,u)=u$. Suppose the observed mediator density in that arm is
\[
 g_1(z,u)=1+c(2z-1)(2u-1),\qquad 0<c<1.
\]
Take a Bernoulli outcome with this controlled mean and an independent response disturbance. The feature $Z$ is uniform, but
$E[Y\mid W=1,Z=z]=1/2+c(2z-1)/6$. Substituting this observed regression in the tilt formula produces a nonzero contrast for every nonzero scalar tilt: the tilted mean of $2Z-1$ is nonzero. The actual reference-fiber mean is $a_1(z)=1/2$, and the correct contrast is zero. This example deliberately drops fiber constancy; it shows the identification work performed by that assumption, rather than contradicting it.
\end{example}

\begin{proposition}[A dimension-free attainable experiment]
Suppose $|Y|\le B$, treatment is randomized with supplied $q_w(X)\ge\epsilon$, and, conditional on $(X,W=w)$, the mediator is drawn from the supplied $g_0$ independently of the response disturbance. Suppose $Z_\theta(x)$ is supplied. For $\Theta=\{\theta:\|\theta\|_1\le R\}\subset\mathbb R^r$, define
\[
 \widehat\psi_w(\theta)=\frac1n\sum_i\frac{1\{W_i=w\}Y_i}{q_w(X_i)}
 \left\{\frac{e^{\theta'h(M_i)}}{Z_\theta(X_i)}-1\right\}.
\]
It is unbiased for every $\theta$. There is a constant $C$ depending only on $B,\epsilon,R$ such that, with probability at least $1-\alpha$,
\[
 \sup_{\theta\in\Theta}|\widehat\psi_w(\theta)-\psi_w(\theta)|
 \le C\left\{\sqrt{\frac{r\log(2+rn)+\log(2/\alpha)}n}+\frac1n\right\}.
\]
In particular the bound has no dependence on the ambient mediator dimension $p$.
\end{proposition}
\begin{proof}
Randomization and iterated expectation give unbiasedness. Since $h_j\in[0,1]$ and $\|\theta\|_1\le R$, the density ratio is at most $e^{2R}$. Its derivative with respect to $\theta_j$ is the ratio times $h_j-E_{g_\theta}h_j$, so the summand is Lipschitz in $\ell^1$ with constant at most $Be^{2R}/\epsilon$, and is bounded in absolute value by $B(e^{2R}+1)/\epsilon$. Cover $\Theta$ by an $\ell^1$ net of radius $1/n$ and cardinality at most $(1+2Rrn)^r$ (round each coordinate on a grid of spacing $1/(rn)$). Hoeffding's inequality and a union bound control all net points. The Lipschitz bounds for both the empirical and population functions extend the bound to every point, adding at most $2Be^{2R}/(\epsilon n)$.
\end{proof}
For a fixed nonconstant tilt, the $n^{-1/2}$ scale cannot be improved even with no covariates: set $\mu_t(z)=1/2+t f(z)$, where $f$ is bounded and has a nonzero integral against the tilt contrast. Two choices $t=\pm c/\sqrt n$ have bounded total relative entropy and target separation of order $n^{-1/2}$. The two-point testing inequality gives a squared-risk lower bound of order $n^{-1}$.

\section*{Remaining gap}
The experiment supplies the reference law and randomizes mediators from it; the original observational problem must estimate the reference feature law and distinguish its fiber average from the observed one. The net bound is not a sharp uniform minimax theorem, and does not analyze unknown smoothness or nuisance estimation. The identity clarifies why low-dimensional features can remove $p$ in favorable designs without establishing the full proposed theory.

\begin{thebibliography}{9}
\bibitem{agenda} C. Cinelli, A. Feller, G. Imbens, E. Kennedy, S. Magliacane, and J. Zubizarreta. \emph{Challenges in Statistics: A Dozen Challenges in Causality and Causal Inference} (2025). \url{https://arxiv.org/html/2508.17099v1}.
\end{thebibliography}
\end{document}
